% ID: SEL-D0016-BC-E05
% ESTADO: cotejada. Segunda pasada visual completa y borrador compilado revisados.
% Fuente: Selectividad/matematicas_aplicadas/matematicas_aplicadas_P_D.pdf
% SHA-256: 11865b83e509d9cc2723c752f7a81f714378e1b651f4a4d010dab4a4576b74c2
% Página: 2; bloque C; ejercicio original 5.

Tres personas se encargan de los cobros de la caja de un supermercado.
El mes pasado, la primera de ellas realizó el 30\% de los cobros, la segunda
el 45\% y la tercera el resto. La dirección del supermercado ha comprobado
que de los cobros realizados por la primera persona, el 1\% son erróneos,
que la segunda cometió errores en el 3\% de los cobros y la tercera en el 2\%.
\begin{enumerate}[label=\alph*)]
  \item \textbf{(1.5 puntos)} Calcule la probabilidad de que un cobro elegido
  al azar haya sido erróneo.
  \item \textbf{(1 punto)} Se elige al azar un cobro correcto.
  \textquestiondown{}Cuál es la probabilidad de que haya sido realizado por
  la segunda persona?
\end{enumerate}
